% Chapter Template

\chapter{Experiments}\singlespacing % Main chapter title

\label{Chapter4} % Change X to a consecutive number; for referencing this chapter elsewhere, use \ref{ChapterX}

\lhead{Chapter 4. \emph{Experiments}} % Change X to a consecutive number; this is for the header on each page - perhaps a shortened title

%----------------------------------------------------------------------------------------
%	SECTION 1
%----------------------------------------------------------------------------------------

\section{Task and Dataset: XNLI}

\textbf{The Cross-lingual Natural Language Inference (XNLI) task} is a benchmark designed to evaluate the capability of large language models (LLMs) to perform reasoning across multiple languages. The goal is to classify the relationship between a premise and a hypothesis into one of three categories: \textit{entailment}, \textit{neutral}, or \textit{contradiction}. This task tests the ability of the model to generalize semantic relationships beyond the English language, making it a challenging benchmark for multilingual models.

\textbf{Dataset Overview:} The XNLI dataset contains annotated sentence pairs across 15 languages. In our experiments, we focus on English (\textit{en}) as the source language and Vietnamese (\textit{vi}), Hindi (\textit{hi}) and Urdu (\textit{ur}) as target languages. The dataset for each language can be represented as:
\begin{equation}
	\mathcal{D}_l = \{(x_p^i, x_h^i, y^i) \,|\, x_p^i, x_h^i \in \mathcal{X}_l, \, y^i \in \{0, 1, 2\}\},
\end{equation}
where \(x_p^i\) and \(x_h^i\) denote the premise and hypothesis sentences in language \(l\), and \(y^i\) is the corresponding label indicating the relationship.

\textbf{Tokenization and Preprocessing:} Each sentence pair \((x_p^i, x_h^i)\) is concatenated with a special token to form the input:
\begin{equation}
	x^i = x_p^i + \text{[EOS]} + x_h^i,
\end{equation}
where \([EOS]\) represents the end-of-sequence token. The concatenated input is tokenized into a sequence of tokens \(\{t_1, t_2, \dots, t_T\}\) with a maximum sequence length \(T = 512\).

\textbf{Train and Test Splits:} The dataset is divided into training and testing subsets for each language:
\begin{equation}
	\mathcal{D}_l = \mathcal{D}_l^\text{train} \cup \mathcal{D}_l^\text{test}, \quad \mathcal{D}_l^\text{train} \cap \mathcal{D}_l^\text{test} = \emptyset.
\end{equation}
For the experiments, a fraction \(\mathcal{F} \subset \mathcal{D}_l^\text{train}\) is used to train the model. The size of \(\mathcal{F}\) is determined as:
\begin{equation}
	|\mathcal{F}| = \text{frac} \times |\mathcal{D}_l^\text{train}|,
\end{equation}
where \(\text{frac}\) is the fraction of the training dataset used.

\textbf{Multilingual Challenges:} In the XNLI task, transferring knowledge from English to other languages like Vietnamese and Hindi requires the model to overcome both semantic and structural differences in language. This transfer is tested by evaluating the zero-shot performance on the target languages without directly training on their data. The objective of this task is to fine-tune LLaMA3, a large language model, on the XNLI dataset to improve cross-lingual inference capabilities. The experiments are designed to test the effectiveness of various fine-tuning setups and analyze the role of language-specific neurons in improving performance across multiple languages.

\section{Baseline Experiments}

The baseline experiments establish a foundation for evaluating the performance of different fine-tuning setups on the XNLI task. We include two baselines: zero-shot evaluation of a model fine-tuned on English (\textit{en}) XNLI data and direct fine-tuning on the target language (\textit{vi}) XNLI data. These baselines provide a comparative framework to assess the impact of fine-tuning specific language neurons using LoRA.

\textbf{Zero-Shot Baseline:} The first baseline involves fine-tuning the pretrained model on the English XNLI dataset. Let \(\mathcal{M}_\text{en}\) denote the model fine-tuned on English data. The zero-shot evaluation involves testing \(\mathcal{M}_\text{en}\) on Vietnamese XNLI datasets without further training. The evaluation metric, accuracy, is defined as:
\begin{equation}
	\text{Accuracy} = \frac{\sum_{i=1}^N \mathbb{I}(\hat{y}^i = y^i)}{N},
\end{equation}
where \(N\) is the total number of test examples, \(y^i\) is the ground truth label, \(\hat{y}^i\) is the predicted label, and \(\mathbb{I}(\cdot)\) is the indicator function.

\textbf{Direct Fine-Tuning Baseline:} The second baseline involves directly fine-tuning the pretrained model on the Vietnamese XNLI dataset. Let \(\mathcal{M}_\text{vi}\) denote the model fine-tuned on Vietnamese data. This baseline establishes an upper bound for performance on the Vietnamese XNLI task, as the model is explicitly trained on the target language data. We again use the accuracy as evaluation metric.

\section{Experiment Goals}

The primary objective of this work is to improve the zero-shot performance of the model on the target language, such as Vietnamese (\textit{vi}), without utilizing any target language training data. This constraint is critical to ensure that the model remains agnostic to the specific characteristics of the target language during training, thereby simulating a real-world cross-lingual transfer learning scenario.

\textbf{Baseline Comparison:} In this setup, we consider two baselines:
\begin{itemize}
	\item \textbf{Baseline 1: Zero-Shot Performance on \textit{vi} using English Training Data.} This baseline involves fine-tuning the pretrained model on the English (\textit{en}) XNLI dataset (\(\mathcal{M}_\text{en}\)) and evaluating its zero-shot performance on the Vietnamese (\textit{vi}) test set. This represents the starting point of cross-lingual transfer without any modifications to the model architecture or targeted interventions.
	\item \textbf{Baseline 2: Performance on \textit{vi} using Direct Fine-Tuning on \textit{vi} Data.} This baseline involves directly fine-tuning the pretrained model on the Vietnamese (\textit{vi}) XNLI dataset (\(\mathcal{M}_\text{vi}\)). It provides an upper bound on performance since it leverages target language training data.
\end{itemize}

\textbf{Objective Constraints:} While the ultimate goal is to improve the zero-shot performance on the \textit{vi} test set, it is imperative to operate under the following constraints:
\begin{enumerate}
	\item No target language (\textit{vi}) training data can be used during the fine-tuning process.
	\item Only the English (\textit{en}) XNLI training data is available for model training and intervention.
	\item The interventions must be limited to model surgery techniques, such as fine-tuning specific neurons or modifying activation patterns.
\end{enumerate}

\textbf{Target Metric:} The improvement in zero-shot performance is quantified as:
\begin{equation}
	\Delta_\text{ZS} = \text{Accuracy}(\mathcal{M}_\text{experiment}) - \text{Accuracy}(\mathcal{M}_\text{en}),
\end{equation}
where \(\mathcal{M}_\text{experiment}\) refers to the model modified through various intervention strategies, and \(\mathcal{M}_\text{en}\) refers to the baseline model fine-tuned on \textit{en}.

\textbf{Goal:} The overarching goal is to maximize \(\Delta_\text{ZS}\), thereby surpassing the performance of Baseline 1 (\(\mathcal{M}_\text{en}\)) through innovative techniques such as:
\begin{itemize}
	\item Fine-tuning language-specific neurons identified for the target language (\textit{vi}) using only the \textit{en} training data.
	\item Utilizing statistical activation patterns (\textit{e.g., mean, quantile-based thresholds}) to intervene in neuron activations during inference.
	\item Applying LoRA (Low-Rank Adaptation) to selectively fine-tune specific components such as \textit{qkv} projections and classification heads while freezing other parts of the model.
\end{itemize}

\textbf{Limitations:} By design, the performance of any intervention cannot exceed Baseline 2 (\(\mathcal{M}_\text{vi}\)), as that baseline uses direct training data from the target language. However, our approach aims to close the gap between Baseline 1 (\(\mathcal{M}_\text{en}\)) and Baseline 2 (\(\mathcal{M}_\text{vi}\)) as much as possible through zero-shot strategies.

\section{Targeted Neuron Fine-tuning Experiment}

This section outlines the experimental setups designed to improve the zero-shot performance of the model on the Vietnamese (\textit{vi}) XNLI dataset. Four configurations are explored, each using different strategies to fine-tune the model, while adhering to the constraint of using only English (\textit{en}) training data during fine-tuning. The experiments aim to utilize targeted interventions to improve performance by selectively tuning neurons and applying LoRA-based adaptations.

\textbf{Setup A: Standard Fine-Tuning with LoRA.} In this setup, the pretrained model is fine-tuned on the English (\textit{en}) XNLI dataset using LoRA. The fine-tuning is applied to the query (\(q\)), key (\(k\)), value (\(v\)), and output (\(o\)) projection matrices (\(W_q^i\), \(W_k^i\), \(W_v^i\), \(W_o^i\)) and the classification head. The resulting model is evaluated on the Vietnamese (\textit{vi}) dataset in a zero-shot setting. Additionally, the \textit{vi}-specific activations are substituted by statistical activation patterns such as mean (\(\mu\)), 90\textsuperscript{th} percentile, and others, defined as:
\begin{equation}
	\hat{h}_{j}^{i,\text{vi}} = \text{stat}_{j}^{i,\text{vi}}, \quad \forall \, \text{stat} \in \{\text{mean}, 90p, 95p, 75p, 5p, 10p, 25p, 0\},
\end{equation}
where \(\hat{h}_{j}^{i,\text{vi}}\) represents the substituted activation for vi language for neuron \((i,j)\) in the MLP block.

\textbf{Setup B: Fine-Tuning with LoRA and English Neurons.} This setup builds on Setup A by incorporating fine-tuning of English-specific neurons identified in the MLP block. Let \(\mathcal{N}_\text{en}\) denote the set of English-specific neurons. During fine-tuning, the LoRA modifications (\(\Delta W\)) are applied to both the query-key-value-output projections and the English neurons, ensuring that:
\begin{equation}
	\Delta W_{j}^{i} = 0, \quad \forall j \notin \mathcal{N}_\text{en}.
\end{equation}
This targeted approach allows the model to specialize in English while preserving the generalization capacity of other neurons.

\textbf{Setup C: Fine-Tuning with LoRA and Target Language Neurons.} In this configuration, the focus shifts to fine-tuning target language such as Vietnamese-specific neurons (\(\mathcal{N}_\text{vi}\)). During fine-tuning, the LoRA modifications (\(\Delta W\)) are applied to both the query-key-value-output projections and the Vietnamese neurons, ensuring that:
\begin{equation}
	\Delta W_{j}^{i} = 0, \quad \forall j \notin \mathcal{N}_\text{vi}.
\end{equation}
This ensures that only the Vietnamese-specific neurons are fine-tuned while others remain untouched.

\textbf{Setup D: Fine-Tuning with LoRA and Both English \& Target Language Neurons.} This setup combines the interventions from Setup B and Setup C, targeting both English and Vietnamese neurons. Let \(\mathcal{N}_\text{en,vi} = \mathcal{N}_\text{en} \cup \mathcal{N}_\text{vi}\) represent the union of English and Vietnamese neurons. During fine-tuning, the LoRA modifications (\(\Delta W\)) are applied to both the query-key-value-output projections and both the English \& Vietnamese neurons, ensuring that:
\begin{equation}
	\Delta W_{j}^{i} = 0, \quad \forall j \notin \mathcal{N}_\text{en,vi}.
\end{equation}
This strategy allows the model to simultaneously adapt to English and Vietnamese language-specific representations.

\textbf{Comparative Analysis:} After fine-tuning, the zero-shot performance on Vietnamese (\textit{vi}) XNLI data is measured using the accuracy metric. The performance improvements achieved through the setups are quantified by:
\begin{equation}
	\Delta_\text{ZS} = \text{Accuracy}(\mathcal{M}_\text{setup}) - \text{Accuracy}(\mathcal{M}_\text{en}),
\end{equation}
where \(\mathcal{M}_\text{setup}\) refers to the model fine-tuned using one of the experimental setups. These improvements are compared against the baselines to evaluate the efficacy of the targeted neuron fine-tuning strategies.







%%%%%%%%%%%%%%%%% Changes%%%%%%%%%%%%%%%%%%%%%%%%














































